\documentclass[10pt]{article}
\usepackage[margin=0.78in]{geometry}
\usepackage{amsmath,amssymb,amsthm,mathtools,microtype,graphicx}
\usepackage[hidelinks]{hyperref}

\newtheorem{theorem}{Theorem}
\newtheorem{lemma}{Lemma}
\newtheorem{corollary}{Corollary}
\DeclareMathOperator{\Var}{Var}
\newcommand{\R}{\mathbb R}
\newcommand{\E}{\mathbb E}

\title{Gaussian laws are exactly the universal carriers of\\normalized log-mgf convexity}
\author{Candidate full classification for the question in arXiv:2311.08351}
\date{17 August 2026}

\begin{document}
\maketitle

\begin{center}
\fbox{\parbox{0.93\textwidth}{\textbf{Result.}
Let $X$ be an $\R^n$-valued random vector whose linear functionals have
finite exponential moments of every real order.  The normalized log-moment
generating function
\[
 \lambda\longmapsto \frac1\lambda\log\E e^{\lambda F(X)},
 \qquad \text{with value $\E F(X)$ at $\lambda=0$},
\]
is convex on $\R$ for every convex $F$ for which all positive exponential
moments are finite if and only if $X$ is a possibly shifted, singular, or
deterministic multivariate Gaussian.  In fact, necessity already follows by
testing affine functions.}}
\end{center}

\section{Source question and precise interpretation}

Chen \cite{Chen} proves that if $g$ is a finite-dimensional Gaussian vector
and $F$ is convex, then
\[
                 \Lambda_{F(g)}(\lambda)
       :=\frac1\lambda\log\E e^{\lambda F(g)}
\]
is convex in $\lambda\in\R$, with the continuous value $\E F(g)$ at zero.
After the theorem the source asks whether the conclusions hold for
\emph{more general random vectors}.  The relevant discussion on source PDF
page 7 is reproduced below.

\begin{figure}[h]
\centering
\includegraphics[width=0.91\textwidth]{figures/open_problem_crop.png}
\caption{The extension question in arXiv:2311.08351, source PDF page 7.}
\end{figure}

There are two different extensions one might seek.  A fixed non-Gaussian
scalar random variable can enjoy convexity of its own normalized log-mgf.
The direct vector-valued analogue of Chen's theorem, however, asks for a
\emph{carrier law}: a law of $X$ for which the conclusion holds after every
convex observation $F$.  The theorem below completely classifies those
laws.

Assume throughout that
\[
             \E e^{\langle u,X\rangle}<\infty
             \qquad(u\in\R^n).
\tag{1}\label{eq:linear-exp}
\]
This is the natural condition making every affine test admissible.  For a
finite convex $F:\R^n\to\R$, define
\[
 \Lambda_{F,X}(\lambda)=
 \begin{cases}
  \lambda^{-1}\log\E e^{\lambda F(X)},&\lambda\ne0,\\
  \E F(X),&\lambda=0.
 \end{cases}
\tag{2}\label{eq:normalized}
\]
We quantify only over $F$ such that the positive exponential moments in
\eqref{eq:normalized} are finite.  Negative moments are then automatic:
every finite convex function has an affine lower support, and
\eqref{eq:linear-exp} applies.

\section{Proof intuition}

Linear functions are convex with both signs.  For a direction $v$, reflection
turns its normalized log-mgf into the negative reflection of the original
function.  Convexity for $v$ and $-v$ therefore makes one function both convex
and concave, hence affine.  Multiplying by $\lambda$ shows that every projected
cumulant generating function is quadratic, which is precisely Gaussianity.
The converse is Chen's theorem after representing a possibly singular
Gaussian vector as an affine image of a standard one.

\section{The classification}

For $v\in\R^n$, let
\[
 K_v(\lambda)=\log\E e^{\lambda\langle v,X\rangle},
 \qquad
 \psi_v(\lambda)=
 \begin{cases}K_v(\lambda)/\lambda,&\lambda\ne0,\\
 \E\langle v,X\rangle,&\lambda=0.
 \end{cases}
\]

\begin{lemma}[The sign-pair rigidity mechanism]\label{lem:sign}
If both $\psi_v$ and $\psi_{-v}$ are convex on $\R$, then $\psi_v$ is
affine and $\langle v,X\rangle$ is Gaussian, allowing zero variance.
\end{lemma}

\begin{proof}
The cumulant generating functions satisfy
\[
 K_{-v}(\lambda)=K_v(-\lambda),
 \qquad
 \psi_{-v}(\lambda)=-\psi_v(-\lambda).
\tag{3}\label{eq:reflection}
\]
Convexity of $\psi_v$ makes $\lambda\mapsto\psi_v(-\lambda)$ convex.
Convexity of the negative of this function, supplied by $\psi_{-v}$, makes
it concave.  Hence $\psi_v$ is simultaneously convex and concave and is
therefore affine, say $\psi_v(\lambda)=a_v+b_v\lambda$.  Consequently
\[
                     K_v(\lambda)=a_v\lambda+b_v\lambda^2.
\]
Differentiating at zero gives
$a_v=\E\langle v,X\rangle$ and
$2b_v=\Var(\langle v,X\rangle)\ge0$.  This is exactly the cumulant
generating function of $N(a_v,2b_v)$; uniqueness of moment generating
functions proves the assertion.
\end{proof}

\begin{theorem}[Universal Gaussian convexity characterizes Gaussian
laws]\label{thm:main}
Under \eqref{eq:linear-exp}, the following are equivalent.
\begin{enumerate}
\item For every finite convex $F:\R^n\to\R$ with
      $\E e^{\lambda F(X)}<\infty$ for every $\lambda>0$, the function
      $\Lambda_{F,X}$ is convex on $\R$.
\item The same assertion holds when $F$ is restricted to affine functions.
\item $X$ is multivariate Gaussian, with arbitrary mean and positive
      semidefinite covariance; singular and deterministic laws are allowed.
\end{enumerate}
\end{theorem}

\begin{proof}
The implication (1)$\Rightarrow$(2) is immediate.  Assume (2).  For every
$v\in\R^n$, apply it to the two linear convex functions
$x\mapsto\langle v,x\rangle$ and $x\mapsto-\langle v,x\rangle$.
Lemma \ref{lem:sign} shows that every one-dimensional projection
$\langle v,X\rangle$ is Gaussian.  By the standard projection
characterization (equivalently, by evaluating the characteristic function
on every $v$), $X$ is a multivariate Gaussian.  This proves (2)$\Rightarrow$(3).

Conversely, write a general Gaussian vector as $X=m+Ag$, where $g$ is a
standard Gaussian vector of some finite dimension and $A$ may be singular.
For convex $F$, the composition $z\mapsto F(m+Az)$ is convex.  Chen's theorem
applied to this composition gives convexity of $\Lambda_{F,X}$, proving
(3)$\Rightarrow$(1).  The deterministic case corresponds to $A=0$ and is
included.
\end{proof}

\begin{corollary}[Minimal test class]
Among laws satisfying \eqref{eq:linear-exp}, it is enough to verify the
convexity for the $2n$ coordinate tests $F(x)=x_j$ and $F(x)=-x_j$ together
with their analogues after every linear change of coordinates.  Equivalently,
the paired tests $F(x)=\pm\langle v,x\rangle$ for all $v$ are complete.
\end{corollary}

\section{An explicit obstruction and sharpness of the scope}

Let $X$ be a symmetric Rademacher random variable.  The identity test is
convex, but its normalized log-mgf is
\[
             \Lambda_X(\lambda)=\frac{\log\cosh\lambda}{\lambda},
             \qquad \Lambda_X(0)=0.
\]
It is not convex on $\R$.  Indeed,
\[
 \Lambda_X(\lambda)=\frac\lambda2-\frac{\lambda^3}{12}
                     +\frac{\lambda^5}{45}+O(\lambda^7),
\]
and directly $\Lambda_X''(1)\approx-0.235652<0$ (while its value at $-1$
has the opposite sign).  Thus even the simplest bounded non-Gaussian law
does not support a blanket extension.

The quantifier over \emph{all} convex observations is essential to the
classification.  Fern\'andez-Unzueta, Melbourne, and Palafox-Castillo
\cite{FMP} subsequently studied scalar variables whose own normalized
log-mgf is convex and proved strict convexity for a broad Ehrhard-like class.
That phenomenon does not contradict Theorem \ref{thm:main}: a non-Gaussian
variable may pass one observation, while the paired affine family forces a
universal carrier law to be Gaussian.

\section{Verification, novelty, and verdict}

The included symbolic script checks the Gaussian affine formula, the
Rademacher series, and the signs of its second derivative at $\pm1$.  These
checks are ancillary; the classification proof is the two-line reflection
identity \eqref{eq:reflection} followed by the projection characterization.

A bounded novelty search on 17 August 2026 covered the run's four cheap
indexes and web searches for the source title, its extension phrase, and
combinations of universal convex tests with Gaussian characterization.  The
most relevant later paper was \cite{FMP}; it develops the scalar
fixed-variable direction and a strictness theorem, but does not state the
universal carrier classification above.  No prior statement of this exact
iff result was found.  Novelty is plausible, not certified.

This packet gives a full answer to the natural universal interpretation of
the source's question.  It does not classify laws for which convexity holds
for one prescribed nonlinear $F$, nor intermediate restricted families of
tests.  Human review should audit the intended quantifier in the source
question and the mild all-linear-exponential-moments hypothesis.  Subject to
that interpretation, the proof is complete.

\begin{thebibliography}{9}
\bibitem{Chen}
W.-K. Chen,
\emph{A Gaussian Convexity for Logarithmic Moment Generating Functions with
Applications in Spin Glasses}, arXiv:2311.08351v2, 2024.

\bibitem{FMP}
M. Fern\'andez-Unzueta, J. Melbourne, and G. Palafox-Castillo,
\emph{On Convex Functions of Gaussian Variables}, arXiv:2510.06676, 2025.
\end{thebibliography}

\end{document}
